\documentclass[11pt,a4paper]{article}
\usepackage[utf8]{inputenc}
\usepackage[T1]{fontenc}
\usepackage[french]{babel}
\usepackage{amsmath,amssymb}
\usepackage{geometry}
\usepackage{enumitem}
\usepackage{booktabs}
\usepackage{array}
\usepackage{tabularx}
\usepackage[table]{xcolor}
\usepackage{tcolorbox}
\usepackage{tikz}
\usepackage{fontawesome5}

\tcbuselibrary{skins,breakable}

\geometry{margin=1.8cm}

% Couleurs
\definecolor{primary}{RGB}{31,78,121}
\definecolor{soft}{RGB}{230,239,247}
\definecolor{amber}{RGB}{178,91,0}
\definecolor{amberbg}{RGB}{255,244,224}
\definecolor{amberborder}{RGB}{240,192,112}
\definecolor{greenbg}{RGB}{233,246,234}
\definecolor{greenborder}{RGB}{182,224,184}
\definecolor{redbg}{RGB}{253,236,236}
\definecolor{redborder}{RGB}{242,179,171}
\definecolor{green}{RGB}{46,125,50}
\definecolor{red}{RGB}{211,47,47}

% Boîtes
\newtcolorbox{idee}[1][]{
  colback=soft, colframe=primary, boxrule=1pt, arc=6pt,
  fonttitle=\bfseries\large, coltitle=white, colbacktitle=primary,
  left=8pt, right=8pt, top=6pt, bottom=6pt,
  breakable, #1
}

\newtcolorbox{retenir}[1][]{
  colback=greenbg, colframe=greenborder, boxrule=1pt, arc=6pt,
  fonttitle=\bfseries\large, coltitle=black, colbacktitle=greenborder,
  left=8pt, right=8pt, top=6pt, bottom=6pt,
  breakable, #1
}

\newtcolorbox{attention}[1][]{
  colback=redbg, colframe=redborder, boxrule=1pt, arc=6pt,
  fonttitle=\bfseries\large, coltitle=black, colbacktitle=redborder,
  left=8pt, right=8pt, top=6pt, bottom=6pt,
  breakable, #1
}

\title{\textcolor{primary}{\textbf{ISO 7870-2 et règles de Nelson — L'essentiel}}\\
\large Synthèse pour étudiants BTS Bioanalyses en Laboratoire de Contrôle}
\author{}
\date{}

\begin{document}

\maketitle

\vspace{-0.5cm}

\begin{idee}[title={\faLightbulb\quad En une phrase}]
L'ISO 7870-2 est le guide international des cartes de contrôle de Shewhart. Les règles de Nelson sont 8 signaux d'alerte qui permettent de détecter une dérive avant qu'elle ne devienne critique.
\end{idee}

\section*{\faBullseye\quad ISO 7870-2 — Carte de contrôle de Shewhart}

\begin{itemize}[leftmargin=1.2cm, itemsep=2pt]
    \item[\textcolor{primary}{\faCheck}] \textbf{Objet :} guide pour l'utilisation et la compréhension des cartes de contrôle de type Shewhart pour la maîtrise statistique des procédés.[reference:0]
    \item[\textcolor{primary}{\faCheck}] \textbf{Portée :} limité aux méthodes de contrôle statistique utilisant le système de cartes de Shewhart.[reference:1]
    \item[\textcolor{primary}{\faCheck}] \textbf{Compléments :} introduction brève aux limites d'alerte, à l'analyse des tendances et à la capabilité du procédé.[reference:2]
    \item[\textcolor{primary}{\faCheck}] \textbf{Versions :} ISO 7870-2:2023 (2\textsuperscript{e} édition) remplace la version 2013.[reference:3]
\end{itemize}

\vspace{0.2cm}

\begin{retenir}[title={\faChartLine\quad Structure d'une carte de Shewhart}]
Trois lignes fondamentales :
\begin{itemize}[leftmargin=*]
    \item \textbf{LC} (ligne centrale) = moyenne du procédé ;
    \item \textbf{LSC} (limite supérieure de contrôle) = LC + \(3\sigma\) ;
    \item \textbf{LIC} (limite inférieure de contrôle) = LC \(-\) \(3\sigma\).
\end{itemize}
Un point hors des limites \(\pm 3\sigma\) signale une cause spéciale.
\end{retenir}

\section*{\faListOl\quad Les 8 règles de Nelson}

Les règles de Nelson (1984) détectent des \textbf{schémas non aléatoires} sur une carte de contrôle. Elles permettent d'identifier une \textbf{dérive} ou un \textbf{dérèglement} du procédé.[reference:4]

\vspace{0.3cm}

\begin{tabularx}{\linewidth}{|c|p{3.2cm}|X|}
\hline
\rowcolor{soft}
\textbf{N°} & \textbf{Nom} & \textbf{Description} \\ \hline
\textbf{1} & Hors limite & 1 point au-delà de \(\pm 3\sigma\). \textbf{Signal immédiat.} \\ \hline
\textbf{2} & Décalage & 9 points consécutifs du même côté de la ligne centrale. \\ \hline
\textbf{3} & Tendance & 6 points consécutifs en augmentation ou en diminution. \\ \hline
\textbf{4} & Alternance & 14 points consécutifs alternant haut-bas. \\ \hline
\textbf{5} & Saut 2\(\sigma\) & 2 points sur 3 au-delà de \(\pm 2\sigma\) du même côté. \\ \hline
\textbf{6} & Saut 1\(\sigma\) & 4 points sur 5 au-delà de \(\pm 1\sigma\) du même côté. \\ \hline
\textbf{7} & Collage & 15 points consécutifs dans la zone \(\pm 1\sigma\) (procédé trop stable). \\ \hline
\textbf{8} & Bimodal & 8 points consécutifs hors de \(\pm 1\sigma\) des deux côtés. \\ \hline
\end{tabularx}

\vspace{0.3cm}

\begin{attention}[title={\faExclamationTriangle\quad À retenir absolument}]
La \textbf{règle n°1} est la plus importante : \textbf{un seul point hors des limites} \(\pm 3\sigma\) suffit à déclarer le procédé « hors contrôle ». Les règles 3 et 5 détectent les \textbf{dérives progressives} typiques en laboratoire (usure, vieillissement des réactifs).
\end{attention}

\vspace{0.3cm}

\begin{retenir}[title={\faGraduationCap\quad Ce que tu dois savoir}]
\begin{enumerate}[leftmargin=1.2cm, itemsep=2pt]
    \item \textbf{Définir} une carte de Shewhart et ses 3 lignes (LC, LSC, LIC).
    \item \textbf{Citer} les règles de Nelson et leur rôle (détection de dérive).
    \item \textbf{Appliquer} la règle n°1 : point hors des limites \(\pm 3\sigma\).
    \item \textbf{Distinguer} cause commune et cause spéciale.
    \item \textbf{Relier} à l'ISO 7870-2 (référentiel international).
\end{enumerate}
\end{retenir}

\vspace{0.4cm}

\begin{center}
\textit{Fiche de synthèse — ISO 7870-2 et règles de Nelson — BTS BioALC — À garder dans ton classeur !}
\end{center}

\end{document}